\begin{textsample}{NEG Dim 6 – global\_south\_2025\_05 – Score 1.00 – global\_south\_2025\_05/124235700.txt}  \label{ex:f6_neg_392}
Top UN court wraps week of hearings on humanitarian aid to Gaza THE HAGUE, May 2 : The top United Nations court on Friday wraps a week of hearings on what Israel must do to ensure desperately needed humanitarian aid reaches Palestinians in Gaza and the occupied West Bank.
Last year, the UN General Assembly asked the International Court of Justice to give an advisory opinion on Israel ' s legal obligations after the country effectively banned the UN agency for Palestinian refugees, the main provider of aid to Gaza, from operating.
Experts say the case could have broader ramifications for the United Nations and its missions worldwide.
The hearings are taking place as the humanitarian aid system in Gaza is nearing collapse and ceasefire efforts remain deadlocked.
Israel has blocked the entry of food, fuel, medicine and other humanitarian supplies since March 2.
It renewed its bombardment on March 18, breaking a ceasefire, and seized large parts of the territory, saying it aims to push Hamas to release more hostages. @ @ @ @ @ @ @ @ @ @ part of its war with Hamas and did not attend the hearing.
The country did provide a 38-page written submission for the court to consider.
What is at stake?
The hearings focused on provision of aid to the Palestinians, but the UN court ' s 15 judges could use their advisory opinion to give legal guidance on the powers of the world body."The court has the opportunity to clarify and address questions about the legal immunities of the United Nations,"Mike Becker, an expert on international human rights law at Trinity College Dublin, told The Associated Press.
Advisory opinions issued by the UN court are described as"nonbinding"as there are no direct penalties attached to ignoring them.
However, the treaty that covers the protections that countries must give to United Nations personnel says that disputes should be resolved through an advisory opinion at the ICJ and the opinion"shall be accepted as decisive by the parties"."The oddity of this particular process,"Becker @ @ @ @ @ @ @ @ @ @ that the opinion is non-binding."What has the ICJ been tasked with answering?
The resolution, sponsored by Norway, seeks the ICJ ' s guidance on"obligations of Israel... in relation to the presence and activities of the United Nations... to ensure and facilitate the unhindered provision of urgently needed supplies essential to the survival of the Palestinian civilian population."The United States, Israel ' s closest ally, voted against it.
Israel ' s ban on the agency, known as UNRWA, which provides aid to Gaza, came into effect in January.
The organisation has faced increased criticism from Prime Minister Benjamin Netanyahu and his far-right allies, who claim the group is deeply infiltrated by Hamas.
UNRWA rejects that claim."We can not let states pick and choose where the UN is going to do its work.
This advisory opinion is a very important opportunity to reinforce that,"Becker said.
Do these proceedings matter for countries other than Israel?
Whatever the judges decide @ @ @ @ @ @ @ @ @ @, according to Juliette McIntyre, an expert on international law at the University of South Australia."Are these immunities absolute or is there wiggle room?
This is useful for where United Nations personnel are working in other places,"McIntyre told AP.
An authoritative answer from the World Court can have influence beyond judicial proceedings as well."Every time a norm is breached, the norm gets weaker.
The advisory opinion in this case could push the norm back,"said McIntyre.
In separate proceedings last year, the court issued an unprecedented and sweeping condemnation of Israel ' s rule over the occupied Palestinian territories, finding Israel ' s presence unlawful and calling for it to end.
The ICJ said Israel had no right to sovereignty in the territories, was violating international laws against acquiring territory by force and was impeding Palestinians ' right to self-determination.
According to McIntyre, the arguments presented this week reflect the opinion handed down just nine months ago."Now the starting premise is that Israel @ @ @ @ @ @ @ @ @ @.
What did the Palestinians and Israelis say?
On Monday, the Palestinian delegation accused Israel of breaching international law in the occupied territories and applauded the move to bring more proceedings to the court."Our journey with the international institutions, be it Security Council, the General Assembly or the ICJ, is we are building things block on top of another block while we are marching towards the accomplishment of the inalienable rights of the Palestinian people, including our right to self-determination, statehood, and the right of the refugees,"Palestinian UN envoy Riyad Mansour told reporters.
Israel has denied it is in violation of international law and said the proceedings are biased.
Israeli Foreign Minister Gideon Saar hit back at the case during a news conference in Jerusalem on Monday."I accuse UNRWA, I accuse the UN, I accuse the secretary-general and I accuse all those that weaponized international law and its institutions in order to deprive the most attacked country in the world, Israel, of @ @ @ @ @ @ @ @ @ @ said.

% matched lemmas: 
\end{textsample}
